\documentclass[11pt]{amsart}
\usepackage[margin=1.05in]{geometry}
\usepackage{amssymb}
\usepackage{booktabs}
\usepackage{array}
\usepackage[hidelinks]{hyperref}
\usepackage{url}
% allow line breaks after hyphens in long URLs
\Urlmuskip=0mu plus 1mu
\expandafter\def\expandafter\UrlBreaks\expandafter{\UrlBreaks\do\-}

\newtheorem{theorem}{Theorem}[section]
\newtheorem{lemma}[theorem]{Lemma}
\newtheorem{proposition}[theorem]{Proposition}
\newtheorem{corollary}[theorem]{Corollary}
\theoremstyle{definition}
\newtheorem{definition}[theorem]{Definition}
\theoremstyle{remark}
\newtheorem{remark}[theorem]{Remark}
\numberwithin{equation}{section}

\newcommand{\R}{\mathbb{R}}
\newcommand{\C}{\mathbb{C}}
\newcommand{\Z}{\mathbb{Z}}
\newcommand{\T}{\mathbb{T}}
\newcommand{\Tr}{\operatorname{Tr}}
\newcommand{\tr}{\operatorname{tr}}
\newcommand{\spec}{\operatorname{spec}}
\newcommand{\Li}{\operatorname{Li}_2}
\newcommand{\Cl}{\operatorname{Cl}_2}
\newcommand{\one}{\mathbf{1}}
\newcommand{\IME}{I_{\mathrm{ME}}}
\newcommand{\FD}{F_{\mathrm{DKZ}}}
\newcommand{\cD}{\mathcal{D}}
\renewcommand{\Re}{\operatorname{Re}}
\renewcommand{\Im}{\operatorname{Im}}
\setlength{\emergencystretch}{3em}
\raggedbottom

\begin{document}

\title[Supplementary Information]{Supplementary Information for\\ ``Fourier measurements are optimal for CGLMP Bell
tests on maximally entangled states in every dimension''}

\author{Ansh Mishra}
\address{Independent researcher, Cumming, GA, USA}
\email{ansh.mishra2025@gmail.com}
\author{Aryan Senthilkumar}
\address{Independent researcher, Johns Creek, GA, USA}
\email{aryan.senthilk@gmail.com}

\subjclass[2020]{81P40; 15A42; 47A63; 42A50; 44A10; 65G30}
\keywords{Bell inequalities, CGLMP inequality, maximally entangled states, rigidity, BMV conjecture, Stahl's
theorem, conjugate function, Stein--Weiss theorem, interval arithmetic}

\begin{abstract}
This Supplementary Information contains the complete proofs of the two analytic theorems of the article
(Theorem~3, the strip inequality for matrices of every size, and Theorem~4, the two-variable
Bessis--Moussa--Villani formula with an explicit positive density), an outline of a second, distribution-free proof of
both, precise statements and proofs of the continuum theorem and of the cell-embedding inequalities, a description of
the computer-assisted certificate CONE$_d$ in its three regimes, the proof of the rigidity theorem, the main theorems
of the Lean~4 formalisation, and the independent verifications and checks. Theorem numbers 1--4 refer to the
article.
\end{abstract}

\maketitle

\tableofcontents

%==================================================================================================
\section{Logical structure}\label{sec:structure}
%==================================================================================================

Throughout, $d\ge2$ is the number of outcomes, $N=4d$, $\tau=\frac1M\Tr$ is the normalized trace on $M_M(\C)$, and
$\IME(d)=\frac{4}{d(d-1)}\sum_{j=1}^{d-1}(d-j)\sec\frac{\pi j}{2d}$. The proof of Theorems~1 and~2 of the article
has the following structure.
\begin{enumerate}
\item[(R)] \emph{Reduction} (Section~\ref{sec:reduction}; proved in earlier work\cite{classical}). Every projective
strategy on $\Phi_D$ gives a family of $N$ projections on some $\C^M$ whose deficit $F-\FD$ is a linear form
$\frac12\langle v,\delta\rangle$ in the ``pair counts'' $\delta$.
\item[(T4)] \emph{Two-variable BMV formula} (Section~\ref{sec:bmv}): an explicit positive density. Analytic.
\item[(T3)] \emph{Strip inequality} (Section~\ref{sec:strip}), deduced from (T4), with an exact defect formula and
its equality case. Analytic.
\item[(C)] \emph{Continuum theorem} for projection-valued step fields on a circle (Section~\ref{sec:cont}), from (T3).
Analytic.
\item[(E)] \emph{Cell inequalities} $\langle u^\ell,\delta\rangle\le0$ for every cell vector $\ell$
(Section~\ref{sec:cells}), from (C). Analytic.
\item[(K)] \emph{CONE$_d$}: $v$ is a non-negative combination of cell vectors $u^\ell$ and trivial directions
(Section~\ref{sec:cone}). Computer-assisted (interval arithmetic) for every $d$.
\item[(A)] (R)+(E)+(K) give Theorem~1 (Section~\ref{sec:cells}).
\item[(B)] The equality case of the chain gives Theorem~2 (Section~\ref{sec:rigidity}).
\end{enumerate}
Nothing else is used. In particular the exact sum-of-squares certificates of earlier work for $3\le d\le20$
\cite{repo} are not used; they give an independent proof of Theorems~1 and~2 for those $d$.

Every item of this list, except (K) for $d\ge21$, has been formalised in Lean~4 (Section~\ref{sec:lean}). For
$2\le d\le20$ the formal proof of Theorems~1 and~2 is complete, including (K), whose certificates are re-evaluated by
the Lean kernel. For $d\ge21$ it is complete except for the single input (K), which is established by the
interval-arithmetic certificates of
Section~\ref{sec:cone}, outside Lean. Theorems~3 and~4 are formalised in full. Thus Theorems~1 and~2 are proved for
every $d$, with a computer-assisted step (interval arithmetic) for $d\ge21$, and fully formalised for $d\le20$.

%==================================================================================================
\section{The CGLMP expression, the clock model and the circle configuration}\label{sec:reduction}
%==================================================================================================

\subsection{CGLMP expression}
With outcomes $A_1,A_2,B_1,B_2\in\Z_d$, the CGLMP expression\cite{CGLMP} is
\begin{multline*}
I_d=\sum_{k=0}^{\lfloor d/2\rfloor-1}\Bigl(1-\frac{2k}{d-1}\Bigr)\Bigl(\bigl[P(A_1=B_1+k)+P(B_1=A_2+k+1)\\
+P(A_2=B_2+k)+P(B_2=A_1+k)\bigr]-\bigl[P(A_1=B_1-k-1)+P(B_1=A_2-k)\\
+P(A_2=B_2-k-1)+P(B_2=A_1-k-1)\bigr]\Bigr),
\end{multline*}
where $P(A_x=B_y+k)$ is the probability that $A_x-B_y\equiv k\pmod d$. Let $m(t)\in\{0,\dots,d-1\}$ be the residue of
$t$ modulo $d$, relabel $X_1=A_2$, $Y_1=B_2$, $X_2=A_1$, $Y_2=B_1$, and put
$S=\mathbb{E}\,m(X_1-Y_1)+\mathbb{E}\,m(Y_1-X_2)+\mathbb{E}\,m(X_2-Y_2)+\mathbb{E}\,m(Y_2-X_1-1)$. Then
$I_d=4-\frac{2}{d-1}S$ for every behaviour (ref.~\cite{classical}, Lemma~2.1), and the local bound $I_d\le2$ is
$S\ge d-1$.

\subsection{Maximally entangled strategies}
For projective measurements $\{P^x_a\},\{Q^y_b\}$ on $\C^D\otimes\C^D$ and $\Phi_D=D^{-1/2}\sum_j|jj\rangle$, one has
$P(a,b|x,y)=\tau(P^x_a(Q^y_b)^{\mathsf T})$ with $\tau=\frac1D\Tr$. With $w=e^{2\pi i/d}$, $U_x=\sum_aw^aP^x_a$,
$U'_y=\sum_bw^bQ^y_b$ and $R_1=U_2$, $R_2=(U'_2)^{\mathsf T}$, $R_3=U_1$, $R_4=(U'_1)^{\mathsf T}$ (so $R_i^d=\one$),
the expansion $m(t)=\frac{d-1}{2}+\sum_{n=1}^{d-1}c_nw^{nt}$, $c_n=-1/(1-w^{-n})$, gives
\begin{equation}\label{eq:Strace}
S=2(d-1)+\sum_{n=1}^{d-1}c_n\sum_{i=0}^3\tau(L^{(i)}_n)
\end{equation}
with the links $L^{(0)}_n=R_1^nR_2^{-n}$, $L^{(1)}_n=R_2^nR_3^{-n}$, $L^{(2)}_n=R_3^nR_4^{-n}$ and
$L^{(3)}_n=w^{-n}R_4^nR_1^{-n}$.
Conversely every quadruple of unitaries with $R_i^d=\one$ comes from a projective strategy on $\Phi_D$.

\subsection{The clock model}
Let $z=e^{2\pi i/N}$. For $V_0,\dots,V_{d-1}\in U(M)$ with $V_k^4=\one$ define the covariant strategy
$W=\sum_k|k\rangle\langle k|\otimes z^kV_k$, $R^V_1=X\otimes\one_M$, $R^V_{i+1}=WR^V_iW^\dagger$, with $X$ the cyclic
shift on $\C^d$.

\begin{theorem}[reduction; ref.~\cite{classical}, Theorem~2.3]\label{thm:reduction}
(1) $I_d(R^V)=4F(V)/(d(d-1))$ with $F(V)=\sum_{j<k}\Re[h_{k-j}\tr(V_jV_k^\dagger)]$, $h_m=\sec\psi_m-i\csc\psi_m$,
$\psi_m=\pi m/(2d)$. (2) Every projective strategy $R$ on $\Phi_D$ is unitarily equivalent to a direct summand of some
$R^V$ (one can take $M=4D$) with $I_d(R)=I_d(R^V)$. The family $V$ is unique up to simultaneous conjugation.
(3) The DKZ strategy is, up to the local unitary $G\otimes\bar G$ with $G=W_0^{-2}$, $W_0=\mathrm{diag}(z^k)$, the
covariant strategy with $M=1$, $V\equiv1$; hence $\IME(d)=4\FD/(d(d-1))$ with $\FD=\sum_{m=1}^{d-1}(d-m)\sec\psi_m$.
\end{theorem}

The family $V$ in (2) is obtained by direct-summing the images of $R$ under the cyclic group of order $4d$ generated by
$\rho(R_1,R_2,R_3,R_4)=(R_2,R_3,R_4,wR_1)$ and applying the Stone--von Neumann theorem for $\Z_d$ to the shift of the
summands; see ref.~\cite{classical}, Appendix~A.

\subsection{The circle configuration}
Let $E_k=V_k^\dagger$. Write $x\in\Z_N$ uniquely as $x=k-da$ ($0\le k\le d-1$, $a\in\Z_4$) and let $Q_x$ be the spectral
projection of $E_k$ for the eigenvalue $i^a$. For each $k$, $\{Q_{k-da}\}_{a\in\Z_4}$ is a projective measurement and
$E_k=\sum_ai^aQ_{k-da}$. We call such a family of $N$ projections a \emph{27B configuration}. Put $A_x=Q_x$,
$B_x=Q_{x+d}$ (so $A_xB_x=0$ and $\sum_xA_x=\sum_xB_x=d\,\one$), $\kappa(u)=\cot(\pi u/N)$ for $u\ne0$, $\kappa(0)=0$, and
\[
\langle A,B\rangle_\tau=\sum_{x,y\in\Z_N}\kappa(x-y)\,\tau(A_xB_y),\qquad T(m)=\sum_y\tau(A_{y+m}B_y),\qquad N(m)=T(m)-T(-m).
\]

\begin{lemma}[linear form]\label{lem:linear}
$F(V)=\frac12\langle A,B\rangle_\tau-\frac d2$, $N(d)=d$, $N(2d-m)=N(m)$, and
\[
F(V)-\FD=\tfrac12\langle v,\delta\rangle,\qquad v_m=2\csc\frac{\pi m}{2d},\qquad \delta_m=N(m)-m\qquad(1\le m\le d-1).
\]
Moreover, for $1\le s\le d-1$, $\delta_s+\delta_{d-s}\le0$.
\end{lemma}

\begin{proof}
The first identity is the projection form of the cotangent representation of ref.~\cite{classical} (Lemma~3.1
there, and its operator version in Section~6); its proof holds verbatim for non-commuting $V_k$. For the others, put
$C(n)=\sum_u\tau(Q_uQ_{u+n})$. Then
\[
T(m)=C(m-d),\qquad N(m)=C(m-d)-C(m+d).
\]
Since
$Q_uQ_{u+2d}=0$ (two outcomes of one measurement) and $\sum_u\tau(Q_u)=d$, $N(d)=C(0)-C(2d)=d$, and, by $C(-n)=C(n)$
and $N$-periodicity, $N(2d-m)=C(d-m)-C(3d-m)=C(m-d)-C(m+d)=N(m)$. As $\kappa$ is odd and
$\kappa(2d)=0$, $\langle A,B\rangle_\tau=\sum_{m=1}^{2d-1}\kappa(m)N(m)=d\kappa(d)+\sum_{m=1}^{d-1}(\kappa(m)+\kappa(2d-m))N(m)$,
and $\kappa(m)+\kappa(2d-m)=\cot\frac{\pi m}{4d}+\tan\frac{\pi m}{4d}=2\csc\frac{\pi m}{2d}$. For $V\equiv1$ the
configuration is $A=[0,d)$, $B=[-d,0)$, with $N(m)=m$. For the last claim, let $\pi_{jk}(t)=\sum_b\tau(Q_{j-db}Q_{k-d(b+t)})$
for sites $j<k$, a probability vector on $\Z_4$. Splitting the sum defining $N(m)$ according to whether the site index
wraps around gives $N(m)=\sum_{k-j=m}[\pi_{jk}(1)-\pi_{jk}(3)]+\sum_{k-j=d-m}[\pi_{jk}(0)-\pi_{jk}(2)]$, so
$N(s)+N(d-s)-d$ is a sum over the $d$ pairs at distance $s$ or $d-s$ of $\pi(0)+\pi(1)-\pi(2)-\pi(3)-1\le0$.
\end{proof}

Hence Theorem~1 is equivalent to $\langle v,\delta\rangle\le0$ for every 27B configuration on every $\C^M$.

%==================================================================================================
\section{Theorem 4: the two-variable BMV formula with an explicit density}\label{sec:bmv}
%==================================================================================================

Throughout this section and the next, $B$ is an orthogonal projection of rank $r$ on $\C^M$, $P=\one-B$, $g$ is
Hermitian, $g_{\rm d}=BgB+PgP$, and $\lambda_{\min}\le\lambda_{\max}$ are the extreme eigenvalues of $g$. $\Tr_P$ and
$\Tr_B$ denote traces over the ranges of $P$ and $B$ of the compressions. For $(a,t)\in\C^2$,
\[
\cD(a,t)=\Tr e^{ag-tP}-e^{-t}\Tr_Pe^{aPgP}-\Tr_Be^{aBgB}.
\]
For $s\in\R$ and $0<\tau<1$ let $p_{s,\tau}(x)=\det(g-s-x(P-\tau))$ and $q_{s,\tau}(x)=\det(g_{\rm d}-s-x(P-\tau))$, let
$x_1(s,\tau),\dots,x_M(s,\tau)$ be the roots of $p_{s,\tau}$ (the eigenvalues of $(P-\tau)^{-1}(g-s)$), and
\[
F(s,\tau)=\frac{1}{2\pi}\sum_{i=1}^M|\Im x_i(s,\tau)|\quad(0<\tau<1),\qquad F(s,\tau)=0\ \text{otherwise}.
\]

\begin{theorem}[Theorem~4 of the article, except the defect formula]\label{thm:bmv}
$F$ is non-negative and continuous on $\R\times(0,1)$, vanishes unless $\lambda_{\min}<s<\lambda_{\max}$, satisfies
$F(s,\tau)\le\frac{M}{2\pi}\|g-s\|/\sqrt{\tau(1-\tau)}$ (so $F\in L^1(\R^2)$ with compact support), and for all
$(a,t)\in\C^2$
\begin{equation}\label{eq:bmv}
\cD(a,t)=a^2\int_0^1\!\!\int_\R e^{as-t\tau}F(s,\tau)\,ds\,d\tau .
\end{equation}
\end{theorem}

\begin{lemma}[pencil facts]\label{lem:pencil}
Let $0<\tau<1$ and $s\in\R$.
(a) Every root of $p_{s,\tau}$ satisfies $|x|\le\|g-s\|/\min(\tau,1-\tau)$, and so does every root of $q_{s,\tau}$.
(b) Every non-real root of $p_{s,\tau}$ satisfies $|x|\le\|g-s\|/\sqrt{\tau(1-\tau)}$.
(c) If $s\le\lambda_{\min}$ or $s\ge\lambda_{\max}$, all roots of $p_{s,\tau}$ are real.
(d) All roots of $q_{s,\tau}$ are real.
(e) $p_{s,\tau}$ and $q_{s,\tau}$ have degree $M$, the same leading coefficient $\det(\tau-P)\ne0$ and the same sum of
roots.
(f) $F$ is continuous on $\R\times(0,1)$.
\end{lemma}

\begin{proof}
A root $x$ comes with a unit vector $\psi$ such that $(g-s)\psi=x(P-\tau)\psi$. (a) $P-\tau$ has eigenvalues $1-\tau$
and $-\tau$, so $|x|\min(\tau,1-\tau)\le|x|\,\|(P-\tau)\psi\|=\|(g-s)\psi\|\le\|g-s\|$; likewise for $q$ with
$\|g_{\rm d}-s\|\le\|g-s\|$. (b) $\langle\psi,(g-s)\psi\rangle=x\langle\psi,(P-\tau)\psi\rangle$ with both inner products
real; if $\Im x\ne0$ then $\langle\psi,(P-\tau)\psi\rangle=0$, i.e. $\|P\psi\|^2=\tau$, $\|B\psi\|^2=1-\tau$, hence
$\|(P-\tau)\psi\|^2=(1-\tau)^2\tau+\tau^2(1-\tau)=\tau(1-\tau)$ and the argument of (a) applies. (c) If $g-s$ is
semidefinite and $\Im x\ne0$, then by (b) $\langle\psi,(g-s)\psi\rangle=0$, so $(g-s)\psi=0$, so $x(P-\tau)\psi=0$ with
$x\ne0$ and $P-\tau$ invertible, a contradiction. (d) $g_{\rm d}-s-x(P-\tau)$ is block diagonal; its determinant
vanishes iff $x=(s-\mu)/\tau$ with $\mu\in\spec(BgB|_{\mathrm{ran}B})$ or $x=(\mu-s)/(1-\tau)$ with
$\mu\in\spec(PgP|_{\mathrm{ran}P})$. (e) $p_{s,\tau}(x)=\det(\tau-P)\det(x-(P-\tau)^{-1}(g-s))$, so the roots sum to
$\Tr[(P-\tau)^{-1}(g-s)]$, and likewise for $q$; since $(P-\tau)^{-1}=P/(1-\tau)-B/\tau$ commutes with $B$,
$\Tr[(P-\tau)^{-1}(g-g_{\rm d})]=\Tr[(P-\tau)^{-1}(BgP+PgB)]=0$. (f) The coefficients of $p_{s,\tau}$ are polynomials in
$(s,\tau)$ and the leading coefficient does not vanish on $(0,1)$, so the roots depend continuously on $(s,\tau)$ as an
unordered $M$-tuple, and $\sum_i|\Im x_i|$ is a continuous symmetric function of it.
\end{proof}

\begin{lemma}[Jensen-type identity]\label{lem:jensen}
Let $p(x)=c\prod_{i=1}^M(x-x_i)$ and $q(x)=c\prod_{i=1}^M(x-y_i)$ with $c\ne0$, all $y_i$ real, and
$\sum_ix_i=\sum_iy_i$. Then $\log|p/q|\in L^1(\R)$ and $\int_\R\log|p(\xi)/q(\xi)|\,d\xi=\pi\sum_i|\Im x_i|$.
\end{lemma}

\begin{proof}
Write $x_i=\alpha_i+i\beta_i$. Then $\log|p/q|=A+C$ with $A(\xi)=\sum_i\frac12\log(1+\beta_i^2/(\xi-\alpha_i)^2)\ge0$
and $C(\xi)=\sum_i[\log|\xi-\alpha_i|-\log|\xi-y_i|]$. Each summand of $A$ is integrable with
$\int_\R\frac12\log(1+b^2/u^2)\,du=\pi|b|$ (substitute $u=|b|v$ and use the antiderivative
$v\log(1+v^{-2})+2\arctan v$). $C$ has only logarithmic singularities, and $C(\xi)=O(\xi^{-2})$ as $|\xi|\to\infty$
because $\sum_i\alpha_i=\sum_iy_i$; so $C\in L^1$. For $L>\max_i(|\alpha_i|,|y_i|)$,
$\int_{-L}^L\log|\xi-a|\,d\xi=m_L(a):=(L-a)\log(L-a)+(L+a)\log(L+a)-2L$, which is even and smooth in $a$ with
$m_L''(a)=2L/(L^2-a^2)$, so $m_L(a)=m_L(0)+a^2/L+O(a^4/L^3)$. Hence
$\int_{-L}^LC=\sum_i(\alpha_i^2-y_i^2)/L+O(L^{-3})\to0$, and $\int_\R C=0$.
\end{proof}

\begin{lemma}[real slices]\label{lem:slices}
For $\xi\in\R$ let $\lambda_j(\xi)$ and $\lambda^0_j(\xi)$ ($j=1,\dots,M$) be the eigenvalues of $g-\xi P$ and of
$g_{\rm d}-\xi P$, and put $U_\xi(w)=\sum_j[(w-\lambda_j(\xi))_+-(w-\lambda^0_j(\xi))_+]$ and
$N_\xi(w)=\sum_j[\one_{w>\lambda_j(\xi)}-\one_{w>\lambda^0_j(\xi)}]$. Then (a) $U_\xi$ is continuous, piecewise linear and
compactly supported, with $U_\xi'=N_\xi$ almost everywhere; (b) $\int_\R e^{aw}U_\xi(w)\,dw=\cD(a,a\xi)/a^2$ for every
complex $a\ne0$; (c) the Hilbert transform $(Hf)(w)=\frac1\pi\,\mathrm{p.v.}\!\int f(u)/(w-u)\,du$ of $N_\xi$ is
$(HN_\xi)(w)=\frac1\pi\log|\det(g-\xi P-w)/\det(g_{\rm d}-\xi P-w)|$ for $w$ outside both spectra.
\end{lemma}

\begin{proof}
(a) Below all eigenvalues $U_\xi=0$; above all of them $U_\xi=\Tr(g_{\rm d}-\xi P)-\Tr(g-\xi P)=0$. (b) $U_\xi''=
\sum_j(\delta_{\lambda_j}-\delta_{\lambda^0_j})$ as distributions and $U_\xi$ has compact support, so two integrations by
parts give $a^{-2}\sum_j(e^{a\lambda_j}-e^{a\lambda^0_j})=a^{-2}[\Tr e^{a(g-\xi P)}-\Tr_Be^{aBgB}-e^{-a\xi}\Tr_Pe^{aPgP}]$.
(c) Pair $\lambda_j$ with $\lambda^0_j$; each difference of indicators is $\pm$ the indicator of the interval between
them, and $\frac1\pi\int_\alpha^\beta du/(w-u)=\frac1\pi\log|(w-\alpha)/(w-\beta)|$.
\end{proof}

\begin{lemma}[limits]\label{lem:limits}
For Hermitian $A$,
\[
\lim_{t\to+\infty}\Tr e^{A-tP}=\Tr_Be^{BAB},\qquad\lim_{t\to-\infty}e^t\,\Tr e^{A-tP}=\Tr_Pe^{PAP}.
\]
\end{lemma}

\begin{proof}
Shifting $A$ by a constant multiplies both sides by the same factor, so let $A\ge\one$. Let $\mu_1\ge\dots\ge\mu_M$ be
the eigenvalues of $H_t=A-tP$ and $\beta_1\ge\dots\ge\beta_r$ those of $BAB|_{\mathrm{ran}B}$. Min-max over subspaces
of $\mathrm{ran}\,B$ gives $\mu_k\ge\beta_k$ ($k\le r$). For $\psi=b+p$ ($b\in\mathrm{ran}B$, $p\in\mathrm{ran}P$) and
$t>\|A\|$, $\langle\psi,H_t\psi\rangle\le\langle b,Ab\rangle+2\|A\|\|b\|\|p\|-(t-\|A\|)\|p\|^2\le\langle b,(A+\epsilon_t)b\rangle$
with $\epsilon_t=\|A\|^2/(t-\|A\|)$, so $\mu_k\le\beta_k+\epsilon_t$ for $k\le r$, and $\mu_{r+1}\le\|A\|-t$ (min-max with
the codimension-$r$ subspace $\mathrm{ran}P$). Hence $\Tr e^{H_t}\to\sum_ke^{\beta_k}$. The second limit is the first
with $B$ and $P$ exchanged, since $e^t\Tr e^{A-tP}=\Tr e^{A+tB}$.
\end{proof}

\begin{proof}[Proof of Theorem~\ref{thm:bmv}]
The properties of $F$ are Lemma~\ref{lem:pencil}. \emph{Step 0.} $\cD$ is entire on $\C^2$. Since $e^{-tP}=B+e^{-t}P$,
$\cD(0,t)=0$ and $\partial_a\cD(0,t)=\Tr(ge^{-tP})-e^{-t}\Tr(PgP)-\Tr(BgB)=0$. Hence
$G(a,t):=\cD(a,t)/a^2=\int_0^1(1-u)\partial_a^2\cD(ua,t)\,du$ is entire.

\emph{Step 1 (a distribution).} Put $\hat G(\zeta)=G(-i\zeta_1,i\zeta_2)$, so that $e^{-i(\zeta_1s+\zeta_2\tau)}=
e^{as-t\tau}$ for $a=-i\zeta_1$, $t=i\zeta_2$. Let $K=[\lambda_{\min},\lambda_{\max}]\times[0,1]$ and
$H_K(\eta)=\sup_{(s,\tau)\in K}(s\eta_1+\tau\eta_2)$. For complex $a,t$ the Hermitian part of $ag-tP$ is
$\eta_1g+\eta_2P$ with $\eta=\Im\zeta$, so $\|e^{ag-tP}\|\le e^{\lambda_{\max}(\eta_1g+\eta_2P)}\le e^{H_K(\eta)}$
(the pair $(\langle\psi,g\psi\rangle,\langle\psi,P\psi\rangle)$ lies in $K$ for unit $\psi$). The two subtracted terms
are bounded in the same way, so $|\cD|\le2Me^{H_K(\eta)}$. For $|\zeta_1|\ge1$ this bounds $\hat G$; for $|\zeta_1|<1$
the maximum principle on $|\zeta_1|\le2$ gives $|\hat G(\zeta)|\le\frac M2e^{H_K(\eta)+3c_K}$ with
$c_K=\max(|\lambda_{\min}|,|\lambda_{\max}|)$. By the Paley--Wiener--Schwartz theorem (ref.~\cite{Hormander},
Theorem~7.3.1) there is a distribution $T$ supported in $K$ with
\begin{equation}\label{eq:T}
T(e^{as-t\tau})=G(a,t)\qquad\text{for all }(a,t)\in\C^2,
\end{equation}
and $\hat G$ is bounded on $\R^2$.

\emph{Step 2 ($T=F$ on the open strip).} Let $\varphi\in C_c^\infty(\R^2)$ with Fourier transform $\hat\varphi$. Since
$T$ has compact support and bounded transform, $T(\varphi)=(2\pi)^{-2}\int_{\R^2}\hat G(k)\hat\varphi(-k)\,dk$. Change
variables $k=(\sigma,-\sigma\xi)$ (Jacobian $|\sigma|$) and use Fubini:
$T(\varphi)=(2\pi)^{-2}\int d\xi\int d\sigma\,|\sigma|\hat G(\sigma,-\sigma\xi)\hat\varphi(-\sigma,\sigma\xi)$.
By Lemma~\ref{lem:slices}(b), $\hat G(\sigma,-\sigma\xi)=\hat U_\xi(\sigma)$, the Fourier transform of $U_\xi$; and
$\hat\varphi(-\sigma,\sigma\xi)=\hat\psi_\xi(-\sigma)$ with $\psi_\xi(w)=\int\varphi(w+\xi\tau,\tau)\,d\tau$. By
Lemma~\ref{lem:slices}(a), $|\sigma|\hat U_\xi=(-i\,\mathrm{sgn}\,\sigma)\hat N_\xi=\widehat{HN_\xi}$, so by Parseval
the inner integral is $2\pi\int\!\!\int(HN_\xi)(s-\xi\tau)\varphi(s,\tau)\,ds\,d\tau$, and by Lemma~\ref{lem:slices}(c)
with $g-\xi P-(s-\xi\tau)=g-s-\xi(P-\tau)$,
\[
T(\varphi)=\frac{1}{2\pi^2}\int_\R d\xi\int\!\!\int\varphi(s,\tau)\log\Bigl|\frac{p_{s,\tau}(\xi)}{q_{s,\tau}(\xi)}\Bigr|
\,ds\,d\tau .
\]
If $\mathrm{supp}\,\varphi\subset[-S,S]\times[\delta,1-\delta]$, all roots of $p_{s,\tau},q_{s,\tau}$ lie in
$|x|\le R=(\|g\|+S)/\delta$ there (Lemma~\ref{lem:pencil}(a)); the integrand is uniformly integrable for $|\xi|\le2R$
(logarithmic singularities), and for $|\xi|>2R$ it is bounded by $2MR^2/\xi^2$, using $|\log|1-z|+\Re z|\le|z|^2$ for
$|z|\le\frac12$ and the equal root sums. So Fubini applies, and Lemma~\ref{lem:jensen} (with
Lemma~\ref{lem:pencil}(d),(e)) gives $T(\varphi)=\int\!\!\int\varphi F$.

\emph{Step 3 (no mass on the edges).} $T-F$ has compact support contained in the two segments
$[\lambda_{\min},\lambda_{\max}]\times\{0\}$ and $\times\{1\}$. By the structure theorem for distributions supported in
a hyperplane (ref.~\cite{Hormander}, Theorem~2.3.5), $T-F=\sum_{k\le n}[\alpha_k\otimes\delta^{(k)}(\tau)+\beta_k\otimes
\delta^{(k)}(\tau-1)]$ with compactly supported distributions $\alpha_k,\beta_k$, and therefore
\[
G(a,t)-\int\!\!\int e^{as-t\tau}F\,ds\,d\tau=\sum_kt^k[A_k(a)+e^{-t}B_k(a)],\qquad A_k(a)=\alpha_k(e^{as}),\ B_k(a)=\beta_k(e^{as}).
\]
Fix real $a\ne0$. As $t\to+\infty$, $G(a,t)\to0$ by Lemma~\ref{lem:limits} and the integral tends to $0$ by dominated
convergence, so the polynomial $\sum_kA_k(a)t^k$ tends to $0$ and all $A_k(a)=0$. As $t\to-\infty$, multiplying by
$e^t$ and using Lemma~\ref{lem:limits} again gives $B_k(a)=0$. The $A_k,B_k$ are entire and vanish on
$\R\setminus\{0\}$, hence identically, so $\alpha_k=\beta_k=0$ and $T=F$. Equation~\eqref{eq:T} is
equation~\eqref{eq:bmv}.
\end{proof}

\begin{remark}[sum rule]\label{rem:sumrule}
The second-order Duhamel formula gives $G(0,t)=\|BgP\|_F^2(1-e^{-t})/t=\|BgP\|_F^2\int_0^1e^{-t\tau}d\tau$. By
Theorem~\ref{thm:bmv}, $\int_\R F(s,\tau)\,ds=\|BgP\|_F^2$ for every $\tau\in(0,1)$. This is the horizontal Radon
projection, which the proof does not use; it is an independent check (Fig.~4d of the article).
\end{remark}

\begin{remark}[$M=2$]\label{rem:m2}
For $B=\mathrm{diag}(1,0)$ and $g=\bigl(\begin{smallmatrix}\alpha&c\\ \bar c&\beta\end{smallmatrix}\bigr)$ the
discriminant of $p_{s,\tau}$ is $(m(\tau)-s)^2-4\tau(1-\tau)|c|^2$ with $m(\tau)=\tau\beta+(1-\tau)\alpha$, so
$F(s,\tau)=\sqrt{4\tau(1-\tau)|c|^2-(s-m(\tau))^2}_+/(2\pi\tau(1-\tau))$: every slice is a semicircle law of mass
$|c|^2$ centred on the segment from $\alpha$ ($\tau=0$) to $\beta$ ($\tau=1$).
\end{remark}

\begin{remark}[structure]
$(P-\tau)^{-1}(g-s)$ is self-adjoint for the indefinite form $[u,u']=\langle u,(P-\tau)u'\rangle$, which has $M-r$
positive and $r$ negative squares. Hence at most $\min(r,M-r)$ conjugate pairs of non-real roots occur, and for
$\mathrm{rank}\,B=1$, $F=\frac1\pi\Im x_+$ for the single root in the upper half-plane (if any).
\end{remark}

\begin{corollary}[Stahl's theorem for projections with an explicit density]
For Hermitian $A$ and an orthogonal projection $P$, $t\mapsto\Tr e^{A-tP}$ is the Laplace transform of the positive
measure $\Tr_B(e^{BAB})\delta_0+\Tr_P(e^{PAP})\delta_1+\bigl(\int_\R e^sF_A(s,\tau)\,ds\bigr)d\tau$ on $[0,1]$, where
$F_A$ is the density of Theorem~\ref{thm:bmv} for $g=A$.
\end{corollary}

\begin{proof}
Take $a=1$, $g=A$ in equation~\eqref{eq:bmv}.
\end{proof}

\begin{remark}[prior work]
Stahl proved that $t\mapsto\Tr e^{A-tB}$ is the Laplace transform of a positive measure for all Hermitian $A$ and
$B$\cite{Stahl,Eremenko}, settling the conjecture of Bessis, Moussa and Villani\cite{BMV,LS}. Hein\"avaara\cite{Hein}
constructed tracial joint spectral measures by a two-dimensional Fourier inversion whose density is given by the
imaginary parts of eigenvalues of an auxiliary pencil, and derived Stahl's theorem from it. Theorem~\ref{thm:bmv} is a
variant of this mechanism, adapted to a projection, to the pinched subtraction and to the division by $a^2$; the same
proof gives the analogous formula for general Hermitian $B$, with atoms at all eigenvalues of $B$. We found the proof
independently and learned of ref.~\cite{Hein} afterwards. What is specific to the present setting is the positivity of
the pinched, $a^2$-divided object and its consequence, the strip inequality.
\end{remark}

\subsection*{An elementary route}
Steps~1--3 of the proof of Theorem~\ref{thm:bmv} use the Paley--Wiener--Schwartz theorem, Fourier inversion and the
structure of distributions supported on a line. They can be replaced by an argument without distributions, based on the
following identity.

\begin{proposition}[Radon identity]\label{prop:RI}
Let $H$ be Hermitian on $\C^M$ and, for $0<\tau<1$, let $y_1(\tau),\dots,y_M(\tau)$ be the eigenvalues of
$(P-\tau)^{-1}H$. Then $\tau\mapsto\sum_i|\Im y_i(\tau)|$ is continuous on $(0,1)$, bounded by
$M\|H\|/\sqrt{\tau(1-\tau)}$, and
\[
\frac1{2\pi}\int_0^1\sum_{i=1}^M\bigl|\Im y_i(\tau)\bigr|\,d\tau=\Tr H_+-\Tr\bigl(BHB+PHP\bigr)_+ .
\]
\end{proposition}

\emph{Outline of the proof.} Replace $H$ by $H+c$ with $\Im c>0$. An eigenvector $\psi$ of a root $y$ gives
$\Im c=\langle\psi,(P-\tau)\psi\rangle\Im y$, so no root is real, and a homotopy $H\to sH$ shows that exactly
$\mathrm{rank}\,P$ roots lie in the upper half-plane. Since $\det(H+c-y(P-\tau))$ depends on $(\tau,c)$ only through
$(y,c+y\tau)$, a simple root satisfies the complex inviscid Burgers equation $\partial_\tau y=y\,\partial_cy$; summed
over the roots in the upper half-plane with a Cauchy integral, which is valid also at repeated roots, this gives the
conservation law $\partial_\tau\sum\mathrm{Log}\,y=\partial_c\sum y$. Integrating over $\tau\in(0,1)$, the end
$\tau\to1$ contributes nothing after the pinched terms are subtracted, and the end $\tau\to0$ gives, through a Schur
complement, $\log\det(H+c)-\log\det(H_{\rm d}+c)$ with $H_{\rm d}=BHB+PHP$. Hence the $\tau$-integral $\Psi(c)$ of the
sum of the roots in the upper half-plane, minus its pinched counterpart, equals
$\mathcal E(c)=\Tr[(H_{\rm d}+c)\mathrm{Log}(H_{\rm d}+c)]-\Tr[(H+c)\mathrm{Log}(H+c)]$ up to a term $-2\pi ikc+\kappa_0$
with $k\in\Z$ and $\kappa_0\in\C$. As $c$ tends to a real $c_0$, dominated convergence (the integrand is bounded by
$M(\|H\|+|c|)(2+(\tau(1-\tau))^{-1/2})$: the argument of Lemma~\ref{lem:pencil}(b) gives
$|y|\le(\|H\|+|c|)/\sqrt{\tau(1-\tau)}$ for the roots in the upper half-plane when $\tau\le\frac12$ and for those in
the lower half-plane when $\tau\ge\frac12$, and the sum of all roots equals that of the pinched pencil)
turns the imaginary parts into $\pi$ times the two sides of the identity for $H+c_0$, up to $-2\pi kc_0+\Im\kappa_0$.
For $c_0>\|H\|$ both sides vanish, because all roots are then real; hence $k=0$, $\Im\kappa_0=0$, and the identity
holds. The complete proof is in the repository\cite{repoOQP}.

\emph{Consequences.} On the line $s=w+\xi\tau$ the roots of $p_{s,\tau}$ are $\xi+y_i$, where the $y_i$ are those of
Proposition~\ref{prop:RI} for $H=g-\xi P-w$, and $BHB+PHP=g_{\rm d}-\xi P-w$. So the identity says that the integral of
$F$ along this line equals $U_\xi(w)$ (Lemma~\ref{lem:slices}): the Radon projections of $F$ are the real slices. By
Lemma~\ref{lem:slices}(b) and Tonelli's theorem, equation~\eqref{eq:bmv} holds for all real $a\ne0$ and real $t$; both
sides are entire on $\C^2$, so it holds everywhere by the identity theorem, applied in $a$ and then in $t$. In the proof
of Theorem~\ref{thm:strip}, step~(b) then needs only a one-variable uniqueness lemma: a continuous $Z$ with
$|Z(w)|\le Ce^{-c|w|}$ whose two-sided Laplace transform vanishes on an interval is zero (analytic continuation to the
imaginary axis and a Gaussian approximate identity). This route, which has been checked by an independent verification
(Section~\ref{sec:checks}), is the one followed by the Lean formalisation (Section~\ref{sec:lean}); there the
one-variable uniqueness step uses the Fourier inversion theorem of Mathlib. The same argument for an arbitrary Hermitian
matrix in place of the projection $P$ gives a proof of Stahl's theorem\cite{Stahl} without distribution theory or
Fourier analysis; we make no priority claim about this by-product.

%==================================================================================================
\section{Theorem 3: the strip inequality}\label{sec:strip}
%==================================================================================================

Let $S=\{0\le\Re\nu\le1\}$. For $\lambda\in\R$ and $x+iy$ with $0<x<1$ let
$h_\lambda(x+iy)=-\frac1{\pi^2}\Im\Li(e^{\pi(y-\lambda)}e^{-i\pi x})$; the argument $-\pi x\in(-\pi,0)$ never meets the
cut $[1,\infty)$ of $\Li$, so $h_\lambda$ is harmonic. For $0<X<1$ let $K_X(u)=\sin(\pi X)/(2(\cosh\pi u-\cos\pi X))$.

\begin{lemma}[strip Poisson kernel]\label{lem:kernel}
$K_X>0$, $K_X(u)\le C_Xe^{-\pi|u|}$, $K_X(u)=\sum_{n\ge1}\sin(n\pi X)e^{-n\pi|u|}$ for $u\ne0$, and for $|\Re a|<\pi$,
$\int_\R e^{-au}K_X(u)\,du=\sin(a(1-X))/\sin a$.
\end{lemma}

\begin{proof}
The series follows from $\sum_nq^n\sin n\theta=q\sin\theta/(1-2q\cos\theta+q^2)$ with $q=e^{-\pi|u|}$,
$\theta=\pi X$. For the transform take $a=i\kappa$, $\kappa\ne0$ real, and integrate
$f(u)=e^{-i\kappa u}/(\cosh\pi u-\cos\pi X)$ over the boundary of $[-L,L]\times[0,2]$. Since $f(u+2i)=e^{2\kappa}f(u)$,
the contour integral is $(1-e^{2\kappa})\int_\R f$; the poles inside are $u=iX$ and $u=i(2-X)$ with residues
$e^{\kappa X}/(i\pi\sin\pi X)$ and $-e^{\kappa(2-X)}/(i\pi\sin\pi X)$. Hence
$\int e^{-i\kappa u}K_X(u)\,du=\sinh(\kappa(1-X))/\sinh\kappa$. Both sides are analytic on $|\Re a|<\pi$.
\end{proof}

\begin{lemma}\label{lem:poisson}
For $0<X<1$, $h_\lambda(X+iY)=\int_\R(w-\lambda)_+K_X(Y-w)\,dw$; $h_\lambda$ extends continuously to $S$ with
$h_\lambda(iY)=(Y-\lambda)_+$ and $h_\lambda(1+iY)=0$.
\end{lemma}

\begin{proof}
The right side $\Psi$ is harmonic ($K_X(Y-w)=\frac12\Re\cot(\pi(z-iw)/2)$ with $z=X+iY$). For $Y<\lambda$ expand the
kernel by Lemma~\ref{lem:kernel} and integrate termwise:
\[
\Psi=\sum_n\frac{\sin(n\pi X)e^{n\pi(Y-\lambda)}}{(n\pi)^2}=-\frac1{\pi^2}\Im\sum_n\frac{(e^{\pi(Y-\lambda)}e^{-i\pi X})^n}{n^2}=h_\lambda .
\]
Two harmonic functions on the connected strip that agree on an open set agree everywhere. As $X\to0$, $K_X(Y-\cdot)$
is an approximate identity at $Y$; as $X\to1$, $K_X\to0$ with exponential tails.
\end{proof}

For $\nu=x+iy\in S$ let $\omega_\nu$ be the measure $K_x(y-w)\,dw$ ($0<x<1$), $\delta_y$ ($x=0$) or $0$ ($x=1$). Then
$h_\lambda(\nu)=\int(w-\lambda)_+\omega_\nu(dw)$ and $\int e^{aw}\omega_\nu(dw)=e^{ay}\sin(a(1-x))/\sin a$
($|\Re a|<\pi$).

\begin{theorem}[Theorem~3 of the article, with the defect formula of Theorem~4]\label{thm:strip}
For every $\lambda\in\R$,
\[
\sum_{\nu\in\spec(B+ig)}h_\lambda(\nu)-\Tr[(P(g-\lambda)P)_+]=V(\lambda):=\int_0^1\!\!\int_\R K_{1-\tau}(s-\lambda)F(s,\tau)\,ds\,d\tau\ \ge0,
\]
with eigenvalues counted with algebraic multiplicity. Equality holds for one (equivalently, every) $\lambda$ if and only
if $[B,g]=0$.
\end{theorem}

\begin{proof}
The eigenvalues $\nu_k=x_k+iy_k$ of $X=B+ig$ satisfy $x_k\in[0,1]$, since $\Re\langle\psi,X\psi\rangle=\langle\psi,B\psi\rangle$.
Let $\rho=\sum_k\omega_{\nu_k}$ and $\sigma=\sum_{\mu\in\spec(PgP|_{\mathrm{ran}P})}\delta_\mu$; the left side is
$\int(w-\lambda)_+d(\rho-\sigma)(w)$.
(a) For real $0<|a|<\pi$, $\hat\rho(a):=\int e^{aw}\rho(dw)=\sum_ke^{ay_k}\sin(a(1-x_k))/\sin a=
[\Tr e^{ag+iaP}-\Tr e^{ag-iaP}]/(2i\sin a)$, because $e^{ag+iaP}=e^{-ia(X-1)}$ and $e^{ag-iaP}=e^{ia(X^\dagger-1)}$. Also
$\hat\sigma(a)=\Tr_Pe^{aPgP}$. The $B$-terms cancel in $\cD(a,-ia)-\cD(a,ia)=2i\sin a\,[\hat\rho(a)-\hat\sigma(a)]$, and
by Theorem~\ref{thm:bmv} at $t=\mp ia$ this equals $2ia^2\int\!\!\int e^{as}\sin(a\tau)F\,ds\,d\tau$. With
Lemma~\ref{lem:kernel} ($X=1-\tau$) and Tonelli,
\[
\hat\rho(a)-\hat\sigma(a)=a^2\int\!\!\int e^{as}\frac{\sin a\tau}{\sin a}F\,ds\,d\tau=a^2\int_\R e^{aw}V(w)\,dw .
\]
$V$ is finite, continuous and $O(e^{-\pi|w|})$ ($\int K_{1-\tau}=\tau$ and $F\le C/\sqrt{\tau(1-\tau)}$).
(b) Hence $\rho$ and $\sigma$ have equal mass and equal first moment. Put $W(w)=\int(w-u)_+d(\rho-\sigma)(u)$; $W$ is
continuous, $O(e^{-\pi|w|})$ at both ends (by the moment identities), and $W''=\rho-\sigma$. Two integrations by parts
give $\int e^{aw}W=a^{-2}(\hat\rho-\hat\sigma)=\int e^{aw}V$ for real $0<|a|<\pi$; both sides are analytic on
$|\Re a|<\pi$, so $W$ and $V$ have the same Fourier transform and $W=V$.
(c) Since $(u-\lambda)_+=(\lambda-u)_++(u-\lambda)$ and $\rho-\sigma$ has zero mass and first moment, the left side
equals $W(\lambda)=V(\lambda)\ge0$.
(d) $K_{1-\tau}>0$ for $0<\tau<1$ and $F\ge0$ is continuous, so $V(\lambda)=0$ forces $F\equiv0$ on
$\R\times(0,1)$, then $\cD\equiv0$ by Theorem~\ref{thm:bmv}, then $\|BgP\|_F^2=0$ by Remark~\ref{rem:sumrule}, i.e.
$[B,g]=0$. Conversely, if $[B,g]=0$ then $p=q$, $F=0$ and $V=0$.
\end{proof}

\begin{remark}[convex order]
By (b)--(c), Theorem~\ref{thm:strip} for all $\lambda$ says that $\sigma$ (the spectrum of the compression $PgP$) is
dominated by $\rho$ (the eigenvalues of $B+ig$ swept to the left edge by harmonic measure) in convex order:
$\sum_{\mu}f(\mu)\le\int f\,d\rho$ for every convex $f$. Applied to $(P,g)$ it gives the mirror statement for $BgB$ and
the right edge.
\end{remark}

\begin{remark}[a second proof in special cases]
For $\mathrm{rank}\,B\le1$, for $\mathrm{corank}\,B\le3$ and for all $M\le5$, Theorem~\ref{thm:strip} also follows
from a one-dimensional reduction (the compression eigenvalues of a rank-one perturbation are quantiles of a Cauchy-type
distribution) combined with an argument-principle rigidity of the critical points of the defect. Earlier, the cases
$M\le4$ and $\mathrm{corank}\,B\le2$ had been proved by a geometric argument based on Pedoe's inequality for polygons and,
independently for $M\le4$, by a Bellman-function argument. These proofs are documented in the repository\cite{repoOQP}.
\end{remark}

%==================================================================================================
\section{The continuum theorem for step fields}\label{sec:cont}
%==================================================================================================

Let $\T=\R/2\pi\Z$, $dm=d\theta/2\pi$, and let $\Phi_c(\alpha,\beta)=\frac1{2\pi^2}[\Cl(2\pi\alpha)+\Cl(2\pi\beta)+\Cl(2\pi(1-\alpha-\beta))]$,
where $\Cl(\theta)=\sum_{n\ge1}\sin(n\theta)/n^2$ is the Clausen function. For $\alpha=\beta=\frac14$,
$\Phi_c=G/\pi^2$ with $G$ Catalan's constant. We need the theorem only for step fields, for which every step can be
made explicit.

\begin{theorem}[continuum theorem for step fields]\label{thm:cont}
Let $t_0,\dots,t_{n-1}$ be distinct points of $\T$ in cyclic order, $J_j=(t_j,t_{j+1})$, and let $B_j$ be orthogonal
projections and $0\le A_j\le\one-B_j$ on $\C^M$. Let $B(\theta)=B_j$, $A(\theta)=A_j$ on $J_j$, and assume
$\int B\,dm=\beta\one$. Let $g$ be the conjugate function of $B$,
\[
g(\theta)=\frac1\pi\sum_j(B_j-B_{j-1})\log\Bigl|\sin\frac{\theta-t_j}{2}\Bigr|\qquad(\theta\notin\{t_j\}).
\]
Then, with $\alpha=\int\tau(A)\,dm$, $\int\tau(Ag)\,dm\le\Phi_c(\alpha,\beta)$. More precisely, for every $\lambda$,
\begin{equation}\label{eq:chain}
\int\tau(Ag)\,dm\le\lambda\alpha+\int\tau\bigl((P(g-\lambda)P)_+\bigr)dm\le\lambda\alpha+h_\lambda(\beta),\qquad P=\one-B,
\end{equation}
and if $\alpha=\beta=\frac14$ and equality holds throughout at $\lambda^*=\log2/(2\pi)$, then $[B(\theta),g(\theta)]=0$ for
almost every $\theta$.
\end{theorem}

\begin{proof}
\emph{Herglotz function.} $H(z)=\int\frac{e^{i\phi}+z}{e^{i\phi}-z}B(\phi)\,dm(\phi)=\beta\one+\frac i\pi\sum_j(B_j-B_{j-1})
\mathrm{Log}(1-ze^{-it_j})$ is analytic on the disc, $H(0)=\beta\one$, and extends continuously to the closed disc minus
$\{e^{it_j}\}$ with boundary values $B(\theta)+ig(\theta)$. For $r\ge\frac12$,
$\|H(re^{i\theta})\|\le C_0+\frac2\pi\sum_j|\log|\sin\frac{\theta-t_j}2||$, a function in every $L^p$.
\emph{Spectrum in the strip.} $\langle\psi,\Re H(z)\psi\rangle=\int P_z\langle\psi,B\psi\rangle\,dm\in(0,1)$ for unit $\psi$
unless $B(\phi)\psi=0$ a.e. or $B(\phi)\psi=\psi$ a.e.; after removing the common kernel and range of the $B_j$ (on which
the statement is trivial) the numerical range of $H(z)$, hence its spectrum, lies in the open strip.
\emph{Harmonicity.} $\phi_\lambda(w)=\frac{i}{\pi^2}\Li(e^{-i\pi w-\pi\lambda})$ is analytic on the open strip with
$\Re\phi_\lambda=h_\lambda$, so $u_\lambda(z)=\Re\Tr\phi_\lambda(H(z))=\sum_{\nu\in\spec H(z)}h_\lambda(\nu)$ is harmonic,
with $u_\lambda(0)=Mh_\lambda(\beta)$.
\emph{Boundary passage.} $0\le h_\lambda(x+iy)\le|y|+|\lambda|+C_1$, so $u_\lambda(re^{i\theta})$ is dominated by an
integrable function; for $\theta\notin\{t_j\}$ it converges to $u^*_\lambda(\theta)=\sum_{\nu\in\spec(B+ig)(\theta)}h_\lambda(\nu)$.
By the mean value property and dominated convergence, $\int u^*_\lambda\,dm=Mh_\lambda(\beta)$.
\emph{Chain.} For $\theta\notin\{t_j\}$, with $X=P(g-\lambda)P$: $\tau(A(g-\lambda))=\tau(AX)\le\tau(AX_+)\le\tau(X_+)$, and
by Theorem~\ref{thm:strip}, $\tau(X_+)\le\frac1Mu^*_\lambda(\theta)$. Integrating gives~\eqref{eq:chain}.
\emph{Legendre identity.} $\min_\lambda[\lambda\alpha+h_\lambda(\beta)]=\Phi_c(\alpha,\beta)$, attained at
$\lambda=\frac1\pi\log\frac{\sin\pi(\alpha+\beta)}{\sin\pi\alpha}$; for $\alpha=\beta=\frac14$ this is $\lambda^*$.
\emph{Equality.} If both inequalities in~\eqref{eq:chain} are equalities at $\lambda^*$, then the non-negative integrand
$\frac1Mu^*_{\lambda^*}-\tau((P(g-\lambda^*)P)_+)$ vanishes a.e., i.e. Theorem~\ref{thm:strip} holds with equality at
$\lambda^*$ for a.e. $\theta$, and its equality clause gives $[B(\theta),g(\theta)]=0$ a.e.
\end{proof}

\begin{remark}
With only the trace constraint $\int\tau(B)=\beta$, the centre value is $\sum_ih_\lambda(\beta_i)\le Mh_\lambda(\beta)$
by concavity of $x\mapsto h_\lambda(x)$ on $[0,1]$ (the $y$-convexity of $h_\lambda$ and harmonicity), and the same
bound follows. The operator constraint holds in all applications below.
\end{remark}

%==================================================================================================
\section{Cell embeddings and the reduction of Theorem 1 to \texorpdfstring{CONE$_d$}{CONE(d)}}\label{sec:cells}
%==================================================================================================

Let $\Delta_d=\{\ell\in\R^d:\ell_r\ge0,\ \sum_r\ell_r=d\}$; extend $\ell_x=\ell_{x\bmod d}$ to $x\in\Z_N$, put $L_0=0$,
$L_{x+1}=L_x+\ell_x$ (so $L_N=N$) and $I_x=[L_x,L_{x+1})\subset\R/N\Z$. For a 27B configuration let $\tilde A,\tilde B$ be
the step fields equal to $A_x,B_x$ on $I_x$, and
\[
Q^\ell=\int\!\!\int\tau(\tilde A(\xi)\tilde B(\eta))\cot\frac{\pi(\xi-\eta)}{N}\,d\xi\,d\eta=\sum_{x,y}K^\ell(x,y)\tau(A_xB_y),
\]
with $K^\ell(x,y)=\int_{I_x}\int_{I_y}\cot\frac{\pi(\xi-\eta)}N\,d\xi\,d\eta$.
With $G(u)=-\frac{N^2}{2\pi^2}\Cl(2\pi u/N)$ (so $G''(u)=\cot(\pi u/N)$), $K^\ell(x,y)$ is the mixed second difference
of $G$ over the cell endpoints.

\begin{proposition}[cell inequalities]\label{prop:cells}
For every $\ell\in\Delta_d$ and every 27B configuration, $Q^\ell\le N^2\Phi_c(\frac14,\frac14)$, with equality for the
DKZ configuration. Averaging over the $d$ rotations $Q_x\mapsto Q_{x-c}$,
\begin{equation}\label{eq:cellineq}
\frac1d\sum_{c=0}^{d-1}\bigl[Q^\ell(\mathrm{rot}_cQ)-N^2\Phi_c(\tfrac14,\tfrac14)\bigr]=\langle u^\ell,\delta\rangle\le0,\qquad
u^\ell_m=\Delta^2_m\Bigl[\frac1d\sum_{r=0}^{d-1}\hat G\bigl(S_r(m)\bigr)\Bigr],
\end{equation}
where $S_r(m)=\ell_r+\dots+\ell_{r+m-1}$, $\hat G(u)=G(u)+G(2d-u)$, $\Delta^2_m$ is the second difference in $m$ and
$1\le m\le d-1$. Moreover $v_m=\hat G''(m)$.
\end{proposition}

\begin{proof}
Rescale $\R/N\Z$ to $\T$. Since $\ell$ depends only on the residue modulo $d$ and $\{Q_{r+jd}\}_j$ is a projective
measurement for each $r$, $\int\tilde B\,dm=\frac1N\sum_r\ell_r\sum_jQ_{r+d+jd}=\frac dN\one=\frac14\one$, and likewise
$\int\tilde A\,dm=\frac14\one$; also $\tilde A\tilde B=0$. Zero-length cells carry no field. Theorem~\ref{thm:cont}
(with $\alpha=\beta=\frac14$; $g$ is the conjugate function of the step field) gives $Q^\ell/N^2=\int\tau(\tilde Ag)\,dm\le
\Phi_c(\frac14,\frac14)$. For DKZ, the windows $[c,c+d)$ and $[c-d,c)$ occupy adjacent arcs of length $d$, whose
continuum value is $N^2\Phi_c$ (Stein--Weiss extremal). Rotations preserve the class of 27B configurations. The kernel
$K^\ell$ is invariant under $(x,y)\mapsto(x+d,y+d)$, so the rotation average is
$\sum_m\tilde K(m)T(m)$ with $\tilde K(m)=\frac1d\sum_{r<d}K^\ell(r+m,r)$, odd in $m$; folding with
$N(2d-m)=N(m)$, $N(d)=d$ as in Lemma~\ref{lem:linear} gives~\eqref{eq:cellineq} with
$u^\ell_m=\tilde K(m)+\tilde K(2d-m)$. Writing $K^\ell(r+m,r)$ through window sums,
$\tilde K(m)=\Delta^2P(m)$ with $P(k)=\frac1d\sum_rG(S_r(k))$, and $S_r(2d-j)=2d-S_{r+d-j}(j)$ gives the window-sum
form. Finally $\hat G''(u)=\cot\frac{\pi u}{4d}+\tan\frac{\pi u}{4d}=2\csc\frac{\pi u}{2d}$.
\end{proof}

\begin{definition}[CONE$_d$]
CONE$_d$ holds if there are a finite positive Borel measure $\mu$ on $\Delta_d$ and $w_s\ge0$ with
$v=\int_{\Delta_d}u^\ell\,\mu(d\ell)+\sum_{s=1}^{\lfloor d/2\rfloor}w_s(e_s+e_{d-s})$.
\end{definition}

\begin{proof}[Proof of Theorem~1, given CONE$_d$]
By Lemma~\ref{lem:linear} and Proposition~\ref{prop:cells},
$2(F-\FD)=\langle v,\delta\rangle=\int\langle u^\ell,\delta\rangle\mu(d\ell)+\sum_sw_s(\delta_s+\delta_{d-s})\le0$.
By Theorem~\ref{thm:reduction}, $I_d\le\IME(d)$ for every projective strategy on every $\Phi_D$.
\end{proof}

\begin{lemma}[Dirichlet cells reproduce the singular part]\label{lem:dirichlet}
If $\ell=d\cdot\mathrm{Dir}(1,\dots,1)$ (the gaps between $d$ uniform random points on a circle of length $d$), then for
every $\phi\in C^2(0,d]$ with $\phi''(s)s^m$ integrable and $1\le m\le d-1$,
$\Delta^2_m\mathbb{E}\,\phi(S(m))=\int_0^d\phi''(s)b_{m,d}(s/d)\,ds=\frac{d}{d+1}\mathbb{E}\,\phi''(dX)$ with
$X\sim\mathrm{Beta}(m+1,d-m+1)$, $b_{m,d}(x)=\binom dmx^m(1-x)^{d-m}$. In particular for $\phi''(s)=1/s$ the result is
exactly $1/m$.
\end{lemma}

\begin{proof}
$S(m)/d\sim\mathrm{Beta}(m,d-m)$ has density $\frac{d-1}{d}b_{m-1,d-2}(s/d)$, and
$b''_{m,d}=d(d-1)(b_{m-2,d-2}-2b_{m-1,d-2}+b_{m,d-2})$; two integrations by parts, in which the point masses of $S(0)=0$
and $S(d)=d$ cancel the boundary terms, give the formula.
\end{proof}

Lemma~\ref{lem:dirichlet} explains the construction of the certificates: the expectation of $u^\ell$ under uniform
Dirichlet cells is a Bernstein--Durrmeyer smoothing of $\hat G''$ which reproduces the pole $2\csc(\pi u/2d)\sim4d/(\pi u)$
exactly, and the residual comes only from the image poles at distance $\ge d$. It is smooth and of relative order $1/d$,
and it must be corrected by non-exchangeable cell distributions.

%==================================================================================================
\section{\texorpdfstring{CONE$_d$ for every $d$}{CONE(d) for every d} (computer-assisted)}\label{sec:cone}
%==================================================================================================

\begin{theorem}[computer-assisted]\label{thm:cone}
CONE$_d$ holds for every $d\ge2$. For $d\le200$, $\mu$ is a finite positive combination of point masses at rational
cell vectors with all cells $\ge\frac1{16}$, and $w=0$. For $d\ge201$, $\mu$ is the law of a Gaussian-modulated
Dirichlet process, which charges only the open simplex.
\end{theorem}

\subsection{\texorpdfstring{$2\le d\le200$}{2 <= d <= 200}: exact rational cells}
For each $d$, $K$ cell vectors $\ell^{(k)}=n^{(k)}/16$ with integers $n^{(k)}_r\ge1$, $\sum_rn^{(k)}_r=16d$
($3.5d\le K\le6.3d$) are given. All window sums are multiples of $\frac1{16}$, so $\hat G$ is needed only at
$j/16$: $\hat G(j/16)=-\frac{N^2}{2\pi^2}[\Cl(\pi j/K_g)+\Cl(\pi(K_g-j)/K_g)]$, $K_g=32d$. Each value is enclosed in
160-bit interval arithmetic by $\Cl(t)=t-t\log t+\sum_{n\le70}\frac{|B_{2n}|t^{2n+1}}{2n(2n+1)!}+[0,\epsilon_{70}(t)]$
($0<t\le\pi$) with exact Bernoulli numbers and an explicit tail bound. A floating-point linear program gives weights
$\tilde\lambda>0$; the exact solution $\lambda=\tilde\lambda+U^{\mathsf T}y$ of $U\lambda=v$ is enclosed through
$|y|\le|R\,\mathrm{res}|/(1-\|I-RUU^{\mathsf T}\|)$ with an approximate inverse $R$. Every weight is certified positive
($d\cdot\min_k\lambda_k\in[0.0816,0.2572]$; smallest $4.25\times10^{-4}$, at $d=198$), and a complete re-verification of
all 199 certificates in interval arithmetic passed. The enclosures of the Clausen function were compared with an
independent 63-point high-precision evaluation. An independent scan of all 199 certificates confirmed the cell
structure and the positivity of the weights, and independent floating-point re-evaluations of the identity at selected
$d$ (with 50-digit Clausen values and, for $d\le10$, the original cell-pair kernel instead of the window-sum form) agree
to the precision of the stored weights. For $2\le d\le20$ the same certificates are also re-evaluated by the Lean
kernel in exact rational interval arithmetic (Section~\ref{sec:lean}).

\subsection{\texorpdfstring{$201\le d\le2000$}{201 <= d <= 2000}: Gaussian-modulated Dirichlet cells}
\emph{Potentials.} For a cell process put $P(m)=\mathbb{E}\frac1d\sum_r\hat G(S_r(m))$, so that $u=\Delta^2P$, and let
$P_v$ solve $\Delta^2P_v=v$ on $1,\dots,d-1$ with the same boundary values. Write $X_A(m)=\frac12(X(m)-X(d-m))$ and
$X_S(m)=\frac12(X(m)+X(d-m))$.

\emph{Lemma R1 (reduction).} If (i) $P_A=(P_v)_A$ and (ii) $W=(P_v-P)_S$ has $\Delta^2W\ge0$ on $1,\dots,d-1$, then
CONE$_d$ holds (with $w=\Delta^2W$, a non-negative combination of the $e_s+e_{d-s}$).

\emph{Lemma R2 (exact singular part).} $\hat G=\hat G_s+\hat G_r$ with $\hat G_s(u)=\frac{4d}\pi u\log\frac ud$ and
$\hat G_r(u)=d^2g_r(u/d)$, $g_r$ analytic on $|x|<2$. For Dirichlet cells the singular part of $P$ equals that of $P_v$
exactly, $\frac{4d}{\pi}[m\psi(m)+1-\frac md-m\psi(d)]$ ($\psi$ the digamma function), so only the analytic $g_r$
enters condition (i).

\emph{Process.} Let $\Phi:\{0,\dots,d\}\to[0,\infty)$ be symmetric with $\Phi(0)=\Phi(d)=0$ and Fourier coefficients
$a_k(\Phi)=-\frac{2w_k}d\sum_m\Phi(m)\cos\frac{2\pi km}d>0$ ($1\le k\le d/2$), where $w_k=1$ for $k<d/2$ and
$w_{d/2}=\frac12$. Then $c(j)=\frac12\Delta^2\Phi(j)$ is a positive-definite
function on $\Z_d$ with zero sum, and the centred stationary Gaussian vector $\eta$ with covariance $c$ satisfies
$\sum_r\eta_r=0$ and $\mathrm{Var}(\eta_0+\dots+\eta_{m-1})=\Phi(m)$. On $E=\{\max_r|\eta_r|<\frac34\}$ take
$\ell=d\cdot\mathrm{Dir}(1+\eta)$, on $E^c$ the uniform Dirichlet law.

\emph{Lemmas R3, R4 (decoupling).} By the aggregation property of the Dirichlet law,
$P(m)=\mathbb{E}[\mathcal F(m+Z(m));E]+\mathbb{P}(E^c)\mathcal F(m)$ with $Z(m)\sim N(0,\Phi(m))$ and
$\mathcal F(b)=\mathbb{E}\hat G(dX_b)$, $X_b\sim\mathrm{Beta}(b,d-b)$. Up to an error bounded by
$2d\exp(-9/(32\Phi(1)))$ times explicit derivative bounds, condition (i) decouples into one scalar equation
$\mathcal H_m(\Phi(m))=\tau_m$ per $1\le m<d/2$.

\emph{Lemma B (fixed point).} If on a box $Q=\prod_m[\mathrm{lo}_m,\mathrm{hi}_m]$ all $a_k>0$, the coupling errors are
at most $2\epsilon_t$, and $\mathcal H_m(\mathrm{lo}_m)<\tau_m-2\epsilon_t$, $\mathcal H_m(\mathrm{hi}_m)>\tau_m+2\epsilon_t$,
then the Poincar\'e--Miranda theorem\cite{Kulpa} gives $\Phi^\#\in Q$ for which (i) holds exactly. The two sign
conditions follow from a certified root enclosure $[a_m,b_m]$ and a certified lower bound $c$ on $\mathcal H_m'$ on
$[a_m-\rho,b_m+\rho]$ with $c\rho\ge2\epsilon_t$.

\emph{Certificates.} For each $201\le d\le2000$ one run of a single checker certifies all hypotheses of Lemma~B and
condition (ii) of Lemma~R1. The run uses one precision (160-bit interval arithmetic throughout) and one code version,
whose SHA-256 hash is stored in every certificate. It encloses $\mathcal F^{(2j)}$ ($j\le7$) at the integers (polygamma
values through zeta values and harmonic sums, and the regular part through exact truncated Pochhammer products with a
rigorous tail), the targets $\tau_m$, and the decoupled solution $\Phi_*$ (exact order-14 Gaussian expansion and
order-16 Lagrange remainder); it finds root boxes $[a_m,b_m]$ with a certified sign change and fixes one
Poincar\'e--Miranda box $Q\supseteq\prod_m[a_m-r,b_m+r]$, with a radius $r=2^{e}$ chosen from the tail bound. All
remaining conditions are certified on this widened box: all $a_k>0$; a lower bound $c$ for $\mathcal H_m'$ with
$c\,r>2\epsilon_t$; the two sign conditions, also directly at the faces of $Q$; a lower bound larger than $4\epsilon_t$ for
the second differences of the symmetric part of $P_v$ minus the Gaussian average, which gives (ii) for the exact
process; and the tail bound, in logarithms. Factorials, binomial coefficients and Bernoulli numbers enter as
exact integers and rationals, the zeta values are enclosed by the Euler--Maclaurin formula with a rigorous remainder,
and no decision is taken on floating-point numbers. An audit script re-checks every condition from the exact decimal
values stored in the certificates: 1800 of 1800 pass. The smallest certified values over all $d$ are
$d^2\min_ka_k\ge0.06319$ (even $d$; at $d=2000$) and $\ge0.12638$ (odd $d$), $d\min_m\Delta^2W\ge0.660$, a ratio of the
certified derivative bound to the required one of $4.0009$ (at least $1$ is needed), face margins at least
$10^{16.7}\cdot2\epsilon_t$, $\epsilon_t\le10^{-31.16}$ and $\mathbb{P}(E^c)\le4.2\times10^{-37}$. (Earlier runs at
100 bits for $d\le600$ gave lower apparent margins, $0.0285$ and $0.0798$; they are superseded.)

\subsection{\texorpdfstring{$d\ge2001$}{d >= 2001}: analytic proof with interval Taylor models}
Let $\epsilon=1/d$. A smooth model $\Phi_e(b)=\epsilon b^2u_e(b)$ is defined, $u_e$ being the unique root in $(0,1)$ of
a quartic $\sum_{j=1}^4\mathrm{Ch}_j(b)u^j=\hat\tau_e(b)$ with explicit coefficients (Binet-type functions of the digamma
derivatives and the Taylor coefficients of $g_r$, and an Euler--Maclaurin interpolant of $P_v$). Positivity of the
$a_k$ reduces, by a Fej\'er--P\'olya argument for the convex sequence $a(m)=\frac16-f(m)$, to the pointwise bounds
$\Delta^4(\Phi_e-F_A)(1)\ge-\gamma\epsilon^2$ and $\Delta^4(\Phi_e-F_A)(m)\ge\beta\epsilon^3$ ($2\le m\le d/2$), where
$F_A$ is an explicit boundary layer, together with $|\Phi_*-\Phi_e|\le\delta_{\max}$. For instance, $\gamma=3$,
$\beta=\frac14$ and $\delta_{\max}=\epsilon^2/20$ would already give $a_k\ge0.011\,w_k\epsilon^2>0$ for every $d\ge2001$;
this is the sense of the sufficient level $\frac14$ in Fig.~5d of the article. The bounds actually certified, by interval
Taylor models in $(x,\epsilon)$, $x=b/d$, on boxes covering every $d\ge2001$ simultaneously, are much better:
\begin{center}
\small
\begin{tabular}{>{\raggedright\arraybackslash}p{4.9cm}>{\raggedright\arraybackslash}p{5.0cm}>{\raggedright\arraybackslash}p{4.3cm}}
\toprule
quantity & region & certified bound\\
\midrule
normalized fourth difference, outer zone & $0.011\le x\le0.501$, $0\le\epsilon\le1/2001$ & $\ge7.419$\\
normalized fourth difference, inner zone & $x\le0.013$, all positions $b\ge1$ & $\ge6.022$\\
first two stencil positions & $m=1$; $m=2$ & $\ge-2.765192\,\epsilon^2$; $\ge206.884\,\epsilon^3$\\
model error & all $m$ & $|\Phi_*-\Phi_e|\le6.6468\times10^{-4}\epsilon^2$\\
Fourier coefficients & all $k$ & $a_k\ge0.1186944\,w_k\epsilon^2$\\
symmetric slack & all $m$ & $\Delta^2W\ge0.620154\,\epsilon$\\
higher Gaussian orders & all $m$ & $|\mathsf G-\mathsf G_4|\le138.103\,\epsilon^3$, $|\partial_t\mathsf G_4|\le11.5703\,d$\\
Fej\'er--P\'olya constant & & $\min\mathcal A\ge0.1054878997577$\\
tail and fixed-point conditions & all $d\ge2001$ & logarithmic margins $>800$\\
\bottomrule
\end{tabular}
\end{center}
Here $\mathsf G(b,t)$ is the truncated Gaussian average of $\mathcal F(b+\cdot)-\mathcal F(b)$ with variance $t$, $\mathsf G_4$ its
expansion to order $t^4$, and $\mathcal A(\theta)=a(0)+2\sum_{m\ge1}a(m)\cos m\theta$. All constants, including these, are
computed in interval arithmetic. A rigour pass completed the written proofs of the model identities (one of them, the
first-order identity of the inner zone, is now proved in full and was checked symbolically and numerically to
$6\times10^{-25}$) and corrected several code-level enclosure defects, one of which had left a sliver of size about
$10^{-20}$ of the parameter range uncovered at $d=2001$; no conclusion changed. An audit script checks the coverage of
all boxes and recomputes every margin from the logs. Non-rigorous cross-checks (for example $d^2\Delta^4(1)=-2.7514$ and
$d^3\Delta^4(m)\approx10$ at $d=4000$) agree with the certified bounds.

\begin{remark}[coverage]
All 1800 single-run certificates for $201\le d\le2000$, one per $d$ and all from the same code version, pass the audit,
the smallest derivative ratio being $4.0009$; the audit for $d\ge2001$ reports every box and every margin as certified.
A single command re-checks all three regimes (the 199 certificates for $d\le200$ against their verification records,
the 1800 certificates for $201\le d\le2000$, and every box and margin for $d\ge2001$) and reports CONE$_d$ certified for
every $d\ge2$. Fig.~5 of the article is generated directly from these certificate files.
\end{remark}

\begin{remark}[finite form]\label{rem:finite}
For $d\ge201$ the certificates give CONE$_d$ in integral form, with a measure $\mu$ carried by the open simplex. A
finite form, $v=\sum_j\alpha_ju^{\ell_j}$ with $\alpha_j\ge0$ and at least one $\ell_j$ with all cells positive and
$\alpha_j>0$, follows from it: the barycentre of $\mu$ lies in the convex hull of $\{u^\ell:\ell\text{ in the open
simplex}\}$, Carath\'eodory's theorem reduces it to at most $d$ points, and each trivial direction $e_s+e_{d-s}$ is a
positive multiple of the cell vector of a cell sequence supported on two residues. This finite form is the input of the
Lean formalisation (Section~\ref{sec:lean}).
\end{remark}

%==================================================================================================
\section{Proof of Theorem 2 (rigidity)}\label{sec:rigidity}
%==================================================================================================

Let a projective strategy $R$ on $\Phi_D$ attain $\IME(d)$; let $V$ be its reduced family and $Q$ its 27B
configuration on $\C^M$. Then $\langle v,\delta\rangle=0$ (Theorem~\ref{thm:reduction}, Lemma~\ref{lem:linear}).

\emph{Step 1.} In the representation of Theorem~\ref{thm:cone} every term is $\le0$
(Proposition~\ref{prop:cells}, Lemma~\ref{lem:linear}), so $\langle u^\ell,\delta\rangle=0$ for $\mu$-almost every
$\ell$. Since $\mu$ charges $\Delta^+_d=\{\ell:\ell_r>0\ \forall r\}$ (all cells $\ge\frac1{16}$ for $d\le200$; Dirichlet
laws with parameters $\ge\frac14$ for $d>200$), there is $\ell^*\in\Delta_d^+$ with $\langle u^{\ell^*},\delta\rangle=0$.
All terms of the rotation average~\eqref{eq:cellineq} are $\le0$ with mean $0$, so
$Q^{\ell^*}(Q)=N^2\Phi_c(\frac14,\frac14)$. (The uniform cell vector would not suffice: $u^{(1,\dots,1)}$ is not
proportional to $v$.)

\emph{Step 2.} By the equality clause of Theorem~\ref{thm:cont}, $[B(\theta),g(\theta)]=0$ for almost every $\theta$,
where $B$ is the step field of $Q$ for the cells of $\ell^*$.

\begin{lemma}[residue lemma]\label{lem:residue}
Let $t_0,\dots,t_{n-1}\in\R$ be distinct modulo $2\pi$ and $J_j\in M_M(\C)$. If
$\sum_jJ_j\cot\frac{\theta-t_j}2=0$ on a non-empty open interval avoiding $\bigcup_j(t_j+2\pi\Z)$, then all $J_j=0$.
\end{lemma}

\begin{proof}
The function $\sum_jJ_j\cot\frac{z-t_j}2$ is holomorphic on the connected set $\C\setminus\bigcup_j(t_j+2\pi\Z)$ and
vanishes on an interval, hence everywhere; its residue at $t_j$ is $2J_j$.
\end{proof}

\emph{Step 3.} On each open arc $J_x$ the field $g$ is real-analytic and $[B_x,g(\theta)]=0$ a.e., hence identically.
Differentiating, $\sum_j[B_x,B_j-B_{j-1}]\cot\frac{\theta-t_j}2=0$ on $J_x$. All cells of $\ell^*$ have positive length,
so the $t_j$ are distinct, and Lemma~\ref{lem:residue} gives $[B_x,B_j]=[B_x,B_{j-1}]$ for all $j$; so
$[B_x,B_j]=[B_x,B_x]=0$. Hence all $Q_x$ commute. (Cells of zero length would hide the corresponding measurement; this
is why Step 1 needs $\ell^*\in\Delta^+_d$.)

\emph{Step 4.} Since $E_k=\sum_ai^aQ_{k-da}$, the $V_k$ commute. By Theorem~B of ref.~\cite{classical} every joint
eigenvector of the $V_k$ carries a one-step configuration, and by Corollary~C of ref.~\cite{classical} the strategy
$R$, a direct summand of $R^V$, is unitarily equivalent to $R^{(0)}\otimes\one_K$ with $d\cdot\dim K=D$. Translating back
to measurements, a unitary $u:\C^D\to\C^d\otimes\C^K$ with $uR_iu^\dagger=R^{\rm DKZ}_i\otimes\one_K$ gives
$uP^x_au^\dagger=P^{x,\rm DKZ}_a\otimes\one$ and $\bar uQ^y_b\bar u^\dagger=Q^{y,\rm DKZ}_b\otimes\one$, i.e. the local
unitary $u\otimes\bar u$ maps the strategy to $\mathrm{DKZ}\otimes\one_K$ and $\Phi_D$ to $\Phi_d\otimes\Phi_K$. The
$R^{\rm DKZ}_i$ generate $M_d(\C)$, so $u$ is unique up to $\one_d\otimes U(K)$. If $d\nmid D$, compactness of the set
of strategies on $\Phi_D$ shows that the maximum is attained and hence is strictly below $\IME(d)$.
\hfill$\square$

%==================================================================================================
\section{Formal verification in Lean 4}\label{sec:lean}
%==================================================================================================

The argument of this Supplementary Information has been formalised in the Lean~4 theorem prover\cite{Lean} with the
mathematical library Mathlib\cite{Mathlib}, in the versions
\begin{center}
Lean 4.33.1,\qquad Mathlib revision \texttt{0df444a360eaa60ab8c11dca51a86af692955474}.
\end{center}
The formalisation has 96 files and about 33,000 lines; it is in the directory \texttt{lean/} of the
repository\cite{repoOQP}. No file uses \texttt{sorry}, \texttt{admit}, an \texttt{axiom} declaration or
\texttt{native\_decide}, and \texttt{\#print axioms} reports, for every main theorem, only the three standard axioms
\texttt{propext}, \texttt{Classical.choice} and \texttt{Quot.sound}. The main theorems (file \texttt{Main.lean} and the
files it imports) are the following.

\begin{center}
\small
\begin{tabular}{>{\raggedright\arraybackslash}p{5.0cm}>{\raggedright\arraybackslash}p{6.9cm}>{\raggedright\arraybackslash}p{2.6cm}}
\toprule
Lean theorem & statement & assumptions\\
\midrule
\nolinkurl{OQP27.maxEntClause_le_twenty} & for every $2\le d\le20$: Theorems~1 and~2, and the DKZ strategy attains
$\IME(d)$ & none\\
\nolinkurl{OQP27.maxEntClause_all} & Theorems~1 and~2 for every $d\ge2$ & CONE$_d$ for $d\ge21$ only
(\nolinkurl{Hyp_ConeCertPos_large})\\
\nolinkurl{OQP27.stripInequality}, \nolinkurl{OQP27.stripEquality} & Theorem~3 for every matrix size $M$, and its
equality case & none\\
\nolinkurl{OQP27.StripL3b.theorem1_bmv2} & Theorem~4 (two-variable BMV formula with explicit density), for all
$(a,t)\in\C^2$ & none\\
\nolinkurl{OQP27.StripL3b.hyp_RI} & the Radon identity (Proposition~\ref{prop:RI}) & none\\
\nolinkurl{OQP27.Red.hyp_reduction} & the exact reduction: strategy $\to$ clock model $\to$ 27B configuration $\to$
linear form (Section~\ref{sec:reduction}) & none\\
\nolinkurl{OQP27.Cell.continuumCell} & the continuum theorem for step fields: Theorem~3 implies every cell inequality
(Sections~\ref{sec:cont}, \ref{sec:cells}) & none\\
\nolinkurl{OQP27.ClassB.classicalTheoremB_all} & the classical (commuting) theorem for every $d$, including the discrete
Hilbert-transform rearrangement ($N=4d$, $|A|=|B|=d$) & none\\
\nolinkurl{OQP27.coneCertPos_le_twenty} & CONE$_d$ for $2\le d\le20$: the explicit certificates of
Section~\ref{sec:cone} re-evaluated by the Lean kernel in exact rational interval arithmetic & none\\
\nolinkurl{OQP27.ConeCertificate.fullCheck_sound} & soundness of the certificate checker & none\\
\nolinkurl{OQP27.dkz_cglmp} & the DKZ strategy attains $\IME(d)$, for every $d$ & none\\
\bottomrule
\end{tabular}
\end{center}

Hence the maximally entangled clause of Open Quantum Problem~27B is fully machine-checked for $2\le d\le20$. For
$d\ge21$ the formal proof is complete except for the single input CONE$_d$, stated in the finite form of
Remark~\ref{rem:finite}, which is established by the interval-arithmetic certificates of Section~\ref{sec:cone},
outside Lean. The formal proof of Theorems~3 and~4 follows the elementary route of Section~\ref{sec:bmv}: the Radon
identity, then Theorem~4 by Tonelli's theorem and the identity theorem; its one-variable uniqueness step uses the
Fourier inversion theorem of Mathlib. The continuum theorem is formalised for step fields, with explicit Herglotz
functions, as in Section~\ref{sec:cont}.

Theorems~3 and~4 are formalised in full in the following sense: the identity~\eqref{eq:bmv} on $\C^2$, with $F\ge0$
measurable, compactly supported and bounded by $C(\tau(1-\tau))^{-1/2}$, and the strip inequality with its exact defect
formula and its equality case, for every matrix size. A few auxiliary statements are proved on paper only, because the
formal proof does not need them: the continuity of $F$ (Lean uses measurability, and a bound in Frobenius norm), the sum
rule of Remark~\ref{rem:sumrule}, and the dilogarithm formula for $h_\lambda$ (Lean defines $h_\lambda$ as the Poisson integral
of Lemma~\ref{lem:poisson}). Likewise, the formal rigidity statement is the main assertion of Theorem~2 ($d$ divides $D$,
and a unitary $u$ maps the measurements to DKZ$\,\otimes\one$); the uniqueness of $u$ and the corollary for $d\nmid D$
(Step~4 of Section~\ref{sec:rigidity}) are proved on paper.

%==================================================================================================
\section{Independent verification and checks}\label{sec:checks}
%==================================================================================================

\emph{Independent verification.} Each proof was refereed by an independent checker, who re-derived it step by step and
wrote the numerical checks from scratch; the reports are in the repository\cite{repoOQP}.
\begin{enumerate}
\item Theorems~3 and~4, by the route of Sections~\ref{sec:bmv} and~\ref{sec:strip} (Paley--Wiener--Schwartz and
Fourier slices): verdict correct. Equation~\eqref{eq:bmv} was confirmed at 82 complex points $(a,t)$ to
$6.8\times10^{-13}$, and the defect formula of Theorem~\ref{thm:strip} at 51 pairs (instance, $\lambda$) to
$1.3\times10^{-12}$.
\item The Radon identity and the elementary route (Proposition~\ref{prop:RI} and its consequences): verdict correct;
the identities were confirmed to 34--46 digits.
\item The reduction chain, the continuum theorem, the cell inequalities, the logic of CONE$_d$ and the rigidity
argument (Sections~\ref{sec:reduction}, \ref{sec:cont}--\ref{sec:rigidity}): verdict correct.
\item The alternative proof of Theorem~3 for $\mathrm{rank}\,B\le1$ (and hence for $\mathrm{corank}\,B\le3$ and all
$M\le5$): verdict correct (no mathematical error; some expository gaps, none of which changes a
statement).
\end{enumerate}

\emph{Checks.} The following checks are not part of the proofs; they guard against errors in them. All scripts and
logs are in the repository\cite{repoOQP}.

\begin{center}
\small
\begin{tabular}{>{\raggedright\arraybackslash}p{4.3cm}>{\raggedright\arraybackslash}p{9.8cm}}
\toprule
item & check\\
\midrule
Theorem~\ref{thm:bmv} & independent re-implementation: equation~\eqref{eq:bmv} to $6.8\times10^{-13}$ (relative) at 82
complex $(a,t)$ on 13 instances ($M\le6$); back-projection identity (step 2) to $10^{-20}$ at 15 random $(s,\tau)$ with
0, 2 and 4 non-real roots; sum rule (Remark~\ref{rem:sumrule}) to $10^{-13}$; Radon identity, which gives the Radon
slices, to 34--46 digits on 22 instances with $M\le8$\\
Theorem~\ref{thm:strip} & defect formula to $1.3\times10^{-12}$ at 51 (instance, $\lambda$) pairs on 9 instances,
including spectra up to $\pm150$, clustered spectra and $M=6$ (independent code); adversarial searches for violations
of the strip inequality for $M\le8$ (none); second proof for rank $\le1$, corank $\le3$, $M\le5$\\
Lemma~\ref{lem:linear}, Prop.~\ref{prop:cells} & identities checked on random non-commuting configurations for
$3\le d\le12$, $M\le4$, to $10^{-13}$; every rotated cell inequality negative on random configurations\\
Theorem~\ref{thm:cont} & chain~\eqref{eq:chain} term by term on step fields ($d=4$, $M=2,3$); mean-value identity;
$\lambda^*$ and $\Phi_c(\frac14,\frac14)=G/\pi^2$ to 40 digits\\
Theorem~\ref{thm:cone} & separate checker and complete re-verification for $d\le200$, and re-evaluation by the Lean
kernel for $d\le20$; floating-point re-evaluation; audit of all 1800 single-run certificates for $201\le d\le2000$;
non-rigorous prototypes for large $d$; audit of box coverage and margins for $d\ge2001$\\
Rigidity & residue formula and Herglotz boundary values against quadrature; gradient ascent: all runs that reach the
maximum end at commuting families\\
Theorems~1 and~2, $3\le d\le20$ & independent exact sum-of-squares certificates\cite{repo}\\
\bottomrule
\end{tabular}
\end{center}

%==================================================================================================
\section{Details of the numerical illustrations}\label{sec:figs}
%==================================================================================================

All figures are produced by \texttt{make\_figures.py}. Fig.~1b: Riemannian gradient ascent with backtracking on
$U(D)^4$ (the four measurement bases, rank-one projectors for $D=5$ and rank-two projectors for $D=10$), 20 runs each
from Haar-random starting points; the ascent direction for a basis $U$ with blocks $\Pi_a$ and gradient operators
$G_a$ is $\sum_a[X_a,\Pi_a]$, $X_a=U^\dagger G_aU$. Figs.~3 and~4: a random instance ($M=5$, $\mathrm{rank}\,B=2$, fixed
seed); $h_\lambda$ is evaluated by the dilogarithm with principal branch; the breakpoints of $F(\cdot,\tau)$ are the
critical values of the eigenvalue curves $x\mapsto\lambda_j(g-x(P-\tau))$, located by bisection; the integrals are
computed by adaptive Gauss--Kronrod quadrature after the substitutions $s=s_a+(s_b-s_a)\sin^2(\theta/2)$ on each
interval between breakpoints and $\tau=\sin^2(\vartheta/2)$. The defect of the strip inequality and the balayage of
$F$ agree to $1.4\times10^{-14}$, and the $\tau$-marginals equal $\|BgP\|_F^2$ to $6\times10^{-14}$. Fig.~5 reads the
certificate files directly: the LP certificates for $d\le200$, the single-run certificates for $201\le d\le2000$ (the
plotted values are the lower endpoints of the certified enclosures), and the box logs of the final run for $d\ge2001$.

\begin{thebibliography}{99}
\bibitem{BMV} D. Bessis, P. Moussa and M. Villani, Monotonic converging variational approximations to the functional
integrals in quantum statistical mechanics, J. Math. Phys. \textbf{16} (1975), 2318--2325.
\bibitem{CGLMP} D. Collins, N. Gisin, N. Linden, S. Massar and S. Popescu, Bell inequalities for arbitrarily
high-dimensional systems, Phys. Rev. Lett. \textbf{88} (2002), 040404.
\bibitem{Eremenko} A. E. Eremenko, Herbert Stahl's proof of the BMV conjecture, Sb. Math. \textbf{206} (2015), 87--92.
\bibitem{Hein} O. Hein\"avaara, Tracial joint spectral measures, preprint, arXiv:2310.03227 (2023).
\bibitem{Hormander} L. H\"ormander, \emph{The Analysis of Linear Partial Differential Operators I}, 2nd edn,
Springer, 1990.
\bibitem{Kulpa} W. Kulpa, The Poincar\'e--Miranda theorem, Amer. Math. Monthly \textbf{104} (1997), 545--550.
\bibitem{Lean} L. de Moura and S. Ullrich, The Lean 4 theorem prover and programming language, in: CADE 28, Lecture
Notes in Comput. Sci. \textbf{12699}, Springer, 2021, 625--635.
\bibitem{LS} E. H. Lieb and R. Seiringer, Equivalent forms of the Bessis--Moussa--Villani conjecture, J. Stat. Phys.
\textbf{115} (2004), 185--190.
\bibitem{Mathlib} The mathlib Community, The Lean mathematical library, in: Proc. 9th ACM SIGPLAN International
Conference on Certified Programs and Proofs (CPP 2020), ACM, 2020, 367--381.
\bibitem{classical} A. Mishra and A. Senthilkumar, A rearrangement inequality for the discrete Hilbert transform, with
an application to optimal CGLMP measurements, preprint (2026), \url{https://github.com/anshM123/CGLMP}.
\bibitem{repo} A. Mishra and A. Senthilkumar, CGLMP on maximally entangled states: exact optimality, rigidity and
Tsirelson bounds, \url{https://github.com/anshM123/CGLMP} (2026).
\bibitem{repoOQP} A. Mishra and A. Senthilkumar, IQOQI Vienna Open Problem 27B: proofs, CONE$_d$ certificates, Lean~4
formalisation and independent verification, \url{https://github.com/anshM123/anshM123-IQOQI-Vienna-Open-Problem-27}
(2026).
\bibitem{Stahl} H. R. Stahl, Proof of the BMV conjecture, Acta Math. \textbf{211} (2013), 255--290.
\end{thebibliography}

\end{document}
